\section{背景与相关工作}
\label{sec:background}

\subsection{分区中性原子架构}

空间光调制器（Spatial Light Modulator, SLM）在存储区和纠缠区形成静态光阱。存储区采用密集阱位保存原子，并与Rydberg照射区域分离；纠缠区的阱位成对设置，同一门位对中的两个原子接受全局Rydberg脉冲，执行受控Z（controlled-Z, CZ）门\cite{evered2023highfidelity,lin2025zac}。单比特旋转由可寻址Raman光束执行。纠缠区内未参与目标门的原子也会受到全局脉冲激发，产生空闲受激误差，因此跨层驻留需要权衡阱转移损耗与额外受激损失\cite{evered2023highfidelity,lin2025zac}。

声光偏转器（Acousto-Optic Deflector, AOD）通过控制两个正交方向的光束偏转，形成行列位置可调的移动光阱阵列。原子从SLM阱转入AOD阱后，随移动光阱输运至目标位置，再转入目标SLM阱\cite{bluvstein2022coherenttransport,lin2025zac}。图~\ref{fig:architecture-preliminaries}(a)中，$q_0$、$q_1$在纠缠区执行双比特门，$q_2$在存储区接受单比特旋转，$q_3$通过AOD从存储区进入纠缠区。

\begin{figure*}[!t]
  \centering
  \resizebox{\textwidth}{!}{% Shared vector-figure vocabulary for the Chinese IEEE manuscript.
% Every figure in this directory inputs this file, so the manuscript only
% needs the standard TikZ package and the libraries already loaded by
% paper_zh.tex.  Keep the guard: figures may be input more than once while
% editing or when a stand-alone smoke-test document is built.
\ifdefined\GALKFigureStyleLoaded\else
\def\GALKFigureStyleLoaded{1}

% Baseline-derived visual tokens: ICCAD uses a restrained blue/orange device
% palette, while ZAC uses conventional black axes and a small fixed plot set.
\definecolor{galkInk}{HTML}{231F20}
\definecolor{galkMuted}{HTML}{666666}
\definecolor{galkGrid}{HTML}{D6D6D6}
\definecolor{galkPaper}{HTML}{FFFFFF}
% Semantic colors are fixed across every figure: orange=entanglement zone,
% blue=storage zone, black=atom identity, red=invalidity, and blue dashes=future
% prediction.  Plot and search-group colours are deliberately independent of
% the two physical-zone colours.
\definecolor{galkStorage}{HTML}{0064BC}
\definecolor{galkStorageFill}{HTML}{E8F1F7}
\definecolor{galkEntangle}{HTML}{E37323}
\definecolor{galkEntangleFill}{HTML}{F9ECE3}
\definecolor{galkResident}{HTML}{485CB3}
\definecolor{galkResidentFill}{HTML}{EEF0F8}
\definecolor{galkGood}{HTML}{A4AF07}
\definecolor{galkBad}{HTML}{C84E45}
\definecolor{galkLook}{HTML}{0064BC}
\definecolor{galkData}{HTML}{315E8A}
\definecolor{galkTail}{HTML}{B44B27}
\definecolor{galkGroupOne}{HTML}{4C78A8}
\definecolor{galkGroupTwo}{HTML}{7A7A7A}
\definecolor{galkGroupThree}{HTML}{9C755F}

\tikzset{
  % Shared engineering-figure contract.  Physical snapshots must place an
  % orange entanglement band above a blue storage field; only their scale may
  % change.  Labels are aligned at the upper-left of both regions.
  galk/base/.style={
    font=\rmfamily\fontsize{9.0pt}{10.0pt}\selectfont,
    text=galkInk,
    line cap=butt,
    line join=miter,
    >=Stealth
  },
  galk/panel title/.style={
    font=\rmfamily\bfseries\fontsize{9.2pt}{10.2pt}\selectfont,
    text=galkInk,
    anchor=west
  },
  galk/label/.style={
    font=\rmfamily\fontsize{9.0pt}{10.0pt}\selectfont,
    text=galkInk
  },
  galk/small label/.style={
    font=\rmfamily\fontsize{8.7pt}{9.7pt}\selectfont,
    text=galkInk
  },
  galk/zone label/.style={
    font=\rmfamily\bfseries\fontsize{8.7pt}{9.7pt}\selectfont,
    text=galkInk,
    anchor=north west,
    inner sep=.5pt
  },
  galk/entanglement zone/.style={
    draw=galkEntangle,
    fill=galkEntangleFill,
    line width=.40pt
  },
  galk/storage zone/.style={
    draw=galkStorage,
    fill=galkStorageFill,
    line width=.40pt
  },
  galk/site/.style={
    circle,
    draw=galkInk!68,
    fill=white,
    line width=.36pt,
    minimum size=1.35mm,
    inner sep=0pt
  },
  galk/occupied atom/.style={
    circle,
    draw=galkInk,
    fill=galkInk,
    line width=.36pt,
    minimum size=1.55mm,
    inner sep=0pt
  },
  galk/target site/.style={
    circle,
    draw=galkInk!60,
    fill=galkInk!10,
    line width=.34pt,
    minimum size=1.55mm,
    inner sep=0pt
  },
  galk/q label/.style={
    font=\rmfamily\itshape\fontsize{8.5pt}{9.5pt}\selectfont,
    text=galkInk,
    inner sep=.25pt
  },
  galk/rydberg pair/.style={
    draw=galkInk!62,
    densely dashed,
    line width=.36pt
  },
  galk/accepted path/.style={
    -{Stealth[length=.95mm,width=.62mm]},
    draw=galkInk,
    line width=.48pt
  },
  galk/alternative path/.style={
    -{Stealth[length=.95mm,width=.62mm]},
    draw=galkMuted,
    densely dashed,
    line width=.48pt
  },
  galk/future path/.style={
    -{Stealth[length=.95mm,width=.62mm]},
    draw=galkLook,
    densely dashed,
    line width=.48pt
  },
  galk/invalid path/.style={
    -{Stealth[length=.95mm,width=.62mm]},
    draw=galkBad,
    densely dashed,
    line width=.50pt
  }
}
\fi

\begin{tikzpicture}[
  galk/base,x=1cm,y=1cm,
  trap/.style={galk/site},
  atom/.style={galk/occupied atom},
  fixed/.style={circle,draw=galkInk!65,line width=.35pt,fill=galkInk!25,
    minimum size=1.55mm,inner sep=0pt},
  motion/.style={galk/accepted path},
  guide/.style={draw=galkStorage!45,densely dashed,line width=.40pt},
  inactive guide/.style={draw=galkGrid,densely dashed,line width=.35pt},
  panel title/.style={galk/panel title},
  relation/.style={font=\rmfamily\fontsize{9.3pt}{10.3pt}\selectfont},
  case label/.style={font=\rmfamily\bfseries\fontsize{9.2pt}{10.2pt}\selectfont},
  explanatory note/.style={font=\rmfamily\fontsize{9.2pt}{10.2pt}\selectfont},
  legend label/.style={font=\rmfamily\fontsize{9.2pt}{10.2pt}\selectfont},
  galk/q label/.append style={font=\rmfamily\itshape\fontsize{9.0pt}{10.0pt}\selectfont},
  constraint divider/.style={draw=galkInk!45,dash pattern=on 2pt off 2pt,line width=.45pt}
]
  \path[use as bounding box] (0,-.65) rectangle (17.80,7.20);

  % Device abstraction: rectangular zones and regular trap arrays.
  \node[panel title] at (.08,6.97) {(a) Zoned neutral-atom array};
  \draw[galk/entanglement zone] (.48,4.48) rectangle (7.40,6.55);
  \draw[galk/storage zone] (.48,1.02) rectangle (7.40,4.10);
  \node[galk/zone label] at (.66,6.39) {entanglement zone};
  \node[galk/zone label] at (.66,3.94) {storage zone};

  % Five equally spaced Rydberg sites, with two traps per site.
  \foreach \cx in {1.18,2.50,3.82,5.14,6.46} {
    \draw[galk/rydberg pair] (\cx,5.12) ellipse (.35 and .25);
    \node[trap] at (\cx-.14,5.12) {};
    \node[trap] at (\cx+.14,5.12) {};
  }
  \node[atom] (qzero) at (3.68,5.12) {};
  \node[atom] (qone) at (3.96,5.12) {};
  \draw[galkInk,line width=.48pt] (qzero)--(qone);
  \node[galk/q label,anchor=north east] at (3.62,4.81) {$q_0$};
  \node[galk/q label,anchor=north west] at (4.02,4.81) {$q_1$};
  \node[anchor=south] (rydberg) at (3.91,5.66) {Rydberg pulse (2Q)};
  \draw[motion] (rydberg.south)--(qone.north);

  % Eleven columns and five rows in the storage zone.
  \foreach \x in {.92,1.51,2.10,2.69,3.28,3.87,4.46,5.05,5.64,6.23,6.82} {
    \foreach \y in {1.42,1.85,2.28,2.71,3.14} {
      \node[trap] at (\x,\y) {};
    }
  }
  \node[atom] (qtwo) at (2.69,3.14) {};
  \node[galk/q label,anchor=north east] at (2.62,3.03) {$q_2$};
  \node[anchor=south] (raman) at (3.87,3.46) {Raman pulse (1Q)};
  \draw[motion] (raman.south west)--(qtwo.north east);

  % A hollow destination and one arrow denote a single transported atom.
  \node[atom] (qthree) at (5.64,2.28) {};
  \node[galk/q label,anchor=north east] at (5.57,2.16) {$q_3$};
  \node[trap] (qthreeDest) at (6.32,5.12) {};
  \draw[motion] (qthree.north east)
    .. controls (6.02,3.12) and (6.23,4.23) .. (qthreeDest.south);
  \node[galk/q label,anchor=north east] at (6.26,4.81) {$q_3$};
  \node[anchor=east] at (5.86,4.28) {AOD transport};
  \draw[galkInk,line width=.38pt] (5.88,4.28)--(6.14,4.28);

  \draw[-{Stealth[length=.9mm,width=.6mm]},line width=.42pt]
    (.26,.84)--(7.53,.84) node[below left,inner sep=1pt] {$x$};
  \draw[-{Stealth[length=.9mm,width=.6mm]},line width=.42pt]
    (.26,.84)--(.26,6.63);
  \node[anchor=center,inner sep=0pt] at (.12,6.48) {$y$};
  \draw[{Stealth[length=.8mm,width=.5mm]}-{Stealth[length=.8mm,width=.5mm]},
    line width=.38pt] (1.18,4.70)--(2.50,4.70);
  \node[fill=galkEntangleFill,inner sep=.7pt] at (1.84,4.70) {$d_{\rm site}$};
  \node[anchor=west,text=galkMuted,explanatory note]
    at (.54,.62) {Schematic; not to scale};

  % Strict ordering and row/column membership are one relation constraint.
  \node[panel title] at (8.02,6.97) {(b) Row/column relation constraint};
  \node[anchor=west] at (8.02,6.57) {Preserve $<$, $=$ and $>$ on both axes};
  \node[case label] at (9.48,6.20) {Valid};
  \node[case label] at (12.65,6.20) {Order reversal};
  \node[case label] at (15.91,6.20) {Row merging};
  \draw[galkGrid,line width=.38pt] (11.01,3.73)--(11.01,6.30);
  \draw[galkGrid,line width=.38pt] (14.27,3.73)--(14.27,6.30);

  % Valid: the common row and the order of two columns are preserved.
  \draw[guide] (8.20,4.72)--(10.79,4.72);
  \draw[guide] (8.20,5.73)--(10.79,5.73);
  \node[atom] (v0s) at (8.47,4.72) {};
  \node[atom] (v1s) at (9.66,4.72) {};
  \node[trap] (v0t) at (8.97,5.73) {};
  \node[trap] (v1t) at (10.16,5.73) {};
  \draw[motion] (v0s)--(v0t);
  \draw[motion] (v1s)--(v1t);
  \node[galk/q label,anchor=north] at (8.47,4.60) {$q_0$};
  \node[galk/q label,anchor=north] at (9.66,4.60) {$q_1$};
  \node[galk/q label,anchor=south] at (8.97,5.83) {$q_0$};
  \node[galk/q label,anchor=south] at (10.16,5.83) {$q_1$};
  \node[relation] at (9.48,4.14) {$x_0<x_1\;\to\;x'_0<x'_1$};
  \node[relation] at (9.48,3.68) {$y_0=y_1\;\to\;y'_0=y'_1$};

  % Invalid: column order reverses, while the common row remains unchanged.
  \draw[guide] (11.26,4.72)--(14.00,4.72);
  \draw[guide] (11.26,5.73)--(14.00,5.73);
  \node[atom] (o0s) at (11.69,4.72) {};
  \node[atom] (o1s) at (13.54,4.72) {};
  \node[trap] (o0t) at (13.54,5.73) {};
  \node[trap] (o1t) at (11.69,5.73) {};
  \draw[galk/invalid path] (o0s)--(o0t);
  \draw[galk/invalid path] (o1s)--(o1t);
  \node[galk/q label,anchor=north] at (11.69,4.60) {$q_0$};
  \node[galk/q label,anchor=north] at (13.54,4.60) {$q_1$};
  \node[galk/q label,anchor=south] at (13.54,5.83) {$q_0$};
  \node[galk/q label,anchor=south] at (11.69,5.83) {$q_1$};
  \node[relation,text=galkBad] at (12.65,4.14)
    {$x_0<x_1\;\to\;x'_0>x'_1$};
  \node[relation] at (12.65,3.68) {$y_0=y_1\;\to\;y'_0=y'_1$};

  % Invalid without a geometric crossing: distinct source rows merge.
  \draw[guide] (14.51,4.72)--(17.40,4.72);
  \draw[guide] (14.51,5.13)--(17.40,5.13);
  \draw[guide] (14.51,5.73)--(17.40,5.73);
  \node[atom] (m0s) at (14.92,4.72) {};
  \node[atom] (m1s) at (16.11,5.13) {};
  \node[trap] (m0t) at (15.51,5.73) {};
  \node[trap] (m1t) at (16.70,5.73) {};
  \draw[galk/invalid path] (m0s)--(m0t);
  \draw[galk/invalid path] (m1s)--(m1t);
  \node[galk/q label,anchor=north] at (14.92,4.60) {$q_0$};
  \node[galk/q label,anchor=north west] at (16.18,5.05) {$q_1$};
  \node[galk/q label,anchor=south] at (15.51,5.83) {$q_0$};
  \node[galk/q label,anchor=south] at (16.70,5.83) {$q_1$};
  \node[relation] at (15.91,4.14) {$x_0<x_1\;\to\;x'_0<x'_1$};
  \node[relation,text=galkBad] at (15.91,3.68)
    {$y_0<y_1\;\to\;y'_0=y'_1$};

  % Separate the two constraint classes without reducing the diagrams or text.
  \draw[constraint divider] (8.02,3.35)--(17.62,3.35);

  % The second constraint is the Cartesian product of activated beams.
  \begin{scope}[yshift=-.45cm]
  \node[panel title] at (8.02,3.33) {(c) Unintended-intersection constraint};
  \node[case label] at (9.48,2.93) {Simultaneous pickup};
  \node[case label] at (12.65,2.93) {Batch $B_1$: pick up $q_0$};
  \node[case label] at (15.91,2.93) {Batch $B_2$: pick up $q_1$};
  \draw[galkGrid,line width=.38pt] (11.01,1.04)--(11.01,3.02);

  % Identical coordinates; the same stationary atom q2 is shown each time.
  \foreach \dx in {0,3.26,6.52} {
    \begin{scope}[xshift=\dx cm]
      \foreach \x in {8.65,10.00} {
        \draw[inactive guide] (\x,1.12)--(\x,2.61);
      }
      \foreach \y in {1.36,2.36} {
        \draw[inactive guide] (8.32,\y)--(10.33,\y);
      }
      \node[trap] at (10.00,1.36) {};
      \node[fixed] at (8.65,2.36) {};
      \node[galk/q label,anchor=east] at (8.49,2.36) {$q_2$};
      \node[galk/q label,anchor=west] at (10.16,2.36) {$q_1$};
    \end{scope}
  }

  % Two rows and two columns also activate the occupied cross-intersection.
  \foreach \x in {8.65,10.00} {
    \draw[guide] (\x,1.12)--(\x,2.61);
  }
  \foreach \y in {1.36,2.36} {
    \draw[guide] (8.32,\y)--(10.33,\y);
  }
  \node[atom] at (8.65,1.36) {};
  \node[atom] at (10.00,2.36) {};
  \node[fixed] at (8.65,2.36) {};
  \node[galk/q label,anchor=north] at (8.65,1.24) {$q_0$};
  \draw[galkBad,line width=.7pt] (8.51,2.22)--(8.79,2.50);
  \draw[galkBad,line width=.7pt] (8.51,2.50)--(8.79,2.22);
  \node[explanatory note,text=galkBad,anchor=center]
    at (9.48,.67) {Unintended pickup of $q_2$};

  % One row--column pair is active in each of the two serialized pickups.
  \draw[guide] (11.91,1.12)--(11.91,2.61);
  \draw[guide] (11.58,1.36)--(13.59,1.36);
  \node[atom] at (11.91,1.36) {};
  \node[atom] at (13.26,2.36) {};
  \node[galk/q label,anchor=north] at (11.91,1.24) {$q_0$};
  \draw[motion] (14.00,1.85)--(14.50,1.85);
  \draw[guide] (16.52,1.12)--(16.52,2.61);
  \draw[guide] (14.84,2.36)--(16.85,2.36);
  \node[trap] at (15.17,1.36) {};
  \node[atom] at (16.52,2.36) {};
  \node[galk/q label,anchor=north] at (15.17,1.24) {vacated};
  \node[explanatory note,anchor=center] at (14.26,.67)
    {Sequential pickup leaves $q_2$ in the SLM};
  \end{scope}

  % A dedicated legend row does not obscure any physical component.
  \begin{scope}[yshift=-.45cm]
  \draw[galkGrid,line width=.38pt] (.48,.32)--(17.62,.32);
  \node[atom] at (.68,.02) {};
  \node[anchor=west,legend label] at (.88,.02) {atom};
  \node[fixed] at (2.09,.02) {};
  \node[anchor=west,legend label] at (2.29,.02) {stationary SLM atom};
  \node[trap] at (5.55,.02) {};
  \node[anchor=west,legend label] at (5.75,.02) {empty / target trap};
  \draw[guide] (8.84,.02)--(9.37,.02);
  \node[anchor=west,legend label] at (9.47,.02) {AOD beam position};
  \draw[motion] (12.39,.02)--(12.92,.02);
  \node[anchor=west,legend label] at (13.02,.02) {valid};
  \draw[galk/invalid path] (14.59,.02)--(15.12,.02);
  \node[anchor=west,legend label] at (15.22,.02) {invalid};
  \end{scope}
\end{tikzpicture}
}
  \caption{分区中性原子架构及两类AOD输运约束。(a) 原子$q_0$、$q_1$执行纠缠门，$q_2$接受局域单比特旋转，$q_3$跨区输运。(b) 行列相对关系约束同时禁止次序交换与行列合并；路径不相交仍可能不满足并行条件。(c) 联合拾取形成的非目标交点覆盖静止原子$q_2$，分为$B_1$、$B_2$可避免干扰；$B_1$后$q_0$的源阱已腾空，而$q_2$仍占原阱。原子标号与轨迹一一对应；区域尺寸和阱位数量为示意。}
  \label{fig:architecture-preliminaries}
\end{figure*}

同一AOD中的并行输运受到两类约束\cite{stade2024abstractmodel,stade2025routingaware}。第一类是\emph{行列相对关系约束}：活动行列在移动中保持次序，同一行或同一列上的原子也不能分离到不同的行列。设两原子的源坐标为$(x_i,y_i)$、$(x_j,y_j)$，目标坐标为$(x'_i,y'_i)$、$(x'_j,y'_j)$。将两条输运编入同一批次，须对任意$\star\in\{<,=,>\}$满足
\begin{equation}
\begin{aligned}
x_i\star x_j&\ \Longleftrightarrow\ x'_i\star x'_j,\\
y_i\star y_j&\ \Longleftrightarrow\ y'_i\star y'_j.
\end{aligned}
\label{eq:aod-relations}
\end{equation}
该条件同时排除行列次序交换和同行同列关系改变。源端分行而目标端合并到同一行，即使图上的路径没有相交，也不能并行执行，见图~\ref{fig:architecture-preliminaries}(b)。

第二类是\emph{非目标交点约束}。同时激活的AOD行列会在所有交点形成光阱，其中未用于目标原子阱转移的交点称为非目标交点（ghost spot）。阱转移时，这些交点不能覆盖静止原子。图~\ref{fig:architecture-preliminaries}(c)中，同时拾取$q_0$、$q_1$会在$q_2$处产生非目标交点；分成$B_1$、$B_2$两批拾取可避免这一干扰。编译器将相容输运编入同一重排批次，各批次顺序执行，每批时延由最长轨迹决定。

批次之间还需满足SLM阱位的占用依赖。图~\ref{fig:architecture-preliminaries}(c)中，$B_1$拾取$q_0$后将其源阱标为腾空：该阱从$q_0$转入移动光阱起即可被占用，这一先后关系取决于源阱的释放时刻，而非$q_0$整段输运的结束时刻\cite{lin2025zac}。同一面板里，$q_2$始终占据原SLM阱$s$；$B_1$与$B_2$只安排$q_0$、$q_1$的先后拾取，并不使$q_2$离开$s$。因此，若另有原子要以$s$为目标，须先迁移$q_2$或改选目标阱；对调这两批的次序不能释放$s$。

本文采用分区中性原子架构的保真度模型评价编译结果\cite{lin2025zac}：
\begin{equation}
\begin{aligned}
\log F={}&N_1\log f_1+N_2\log f_2
+N_{\mathrm{exc}}\log f_{\mathrm{exc}}\\
&+N_{\mathrm{tran}}\log f_{\mathrm{tran}}
+\sum_{q\in\mathcal Q}\log\left(1-\frac{t_q}{T_2}\right).
\end{aligned}
\label{eq:fidelity}
\end{equation}
其中，$\mathcal Q$为原子集合，$N_1$、$N_2$分别为单、双比特门数，$N_{\mathrm{exc}}$为纠缠区非参与原子的受激次数，$N_{\mathrm{tran}}$为SLM与AOD之间的逐原子阱转移次数；$f_1$、$f_2$、$f_{\mathrm{exc}}$和$f_{\mathrm{tran}}$为相应的单次保真度因子，$t_q$为原子$q$的累计空闲时间，$T_2$为相干时间。门数固定时，编译主要通过减少阱转移、空闲受激和相干损失提高保真度。参照Stade等人的路由评价\cite{stade2025routingaware}，本文统计重排批次与重排时延。重排批次为必须顺序执行的相容轨迹组数，批次越少，路由并行性越高；重排时延为各批次关键轨迹持续时间之和，反映完成原子输运所需的时间。

\subsection{中性原子编译方法}

\begin{figure*}[t]
  \centering
  \resizebox{\textwidth}{!}{% Shared vector-figure vocabulary for the Chinese IEEE manuscript.
% Every figure in this directory inputs this file, so the manuscript only
% needs the standard TikZ package and the libraries already loaded by
% paper_zh.tex.  Keep the guard: figures may be input more than once while
% editing or when a stand-alone smoke-test document is built.
\ifdefined\GALKFigureStyleLoaded\else
\def\GALKFigureStyleLoaded{1}

% Baseline-derived visual tokens: ICCAD uses a restrained blue/orange device
% palette, while ZAC uses conventional black axes and a small fixed plot set.
\definecolor{galkInk}{HTML}{231F20}
\definecolor{galkMuted}{HTML}{666666}
\definecolor{galkGrid}{HTML}{D6D6D6}
\definecolor{galkPaper}{HTML}{FFFFFF}
% Semantic colors are fixed across every figure: orange=entanglement zone,
% blue=storage zone, black=atom identity, red=invalidity, and blue dashes=future
% prediction.  Plot and search-group colours are deliberately independent of
% the two physical-zone colours.
\definecolor{galkStorage}{HTML}{0064BC}
\definecolor{galkStorageFill}{HTML}{E8F1F7}
\definecolor{galkEntangle}{HTML}{E37323}
\definecolor{galkEntangleFill}{HTML}{F9ECE3}
\definecolor{galkResident}{HTML}{485CB3}
\definecolor{galkResidentFill}{HTML}{EEF0F8}
\definecolor{galkGood}{HTML}{A4AF07}
\definecolor{galkBad}{HTML}{C84E45}
\definecolor{galkLook}{HTML}{0064BC}
\definecolor{galkData}{HTML}{315E8A}
\definecolor{galkTail}{HTML}{B44B27}
\definecolor{galkGroupOne}{HTML}{4C78A8}
\definecolor{galkGroupTwo}{HTML}{7A7A7A}
\definecolor{galkGroupThree}{HTML}{9C755F}

\tikzset{
  % Shared engineering-figure contract.  Physical snapshots must place an
  % orange entanglement band above a blue storage field; only their scale may
  % change.  Labels are aligned at the upper-left of both regions.
  galk/base/.style={
    font=\rmfamily\fontsize{9.0pt}{10.0pt}\selectfont,
    text=galkInk,
    line cap=butt,
    line join=miter,
    >=Stealth
  },
  galk/panel title/.style={
    font=\rmfamily\bfseries\fontsize{9.2pt}{10.2pt}\selectfont,
    text=galkInk,
    anchor=west
  },
  galk/label/.style={
    font=\rmfamily\fontsize{9.0pt}{10.0pt}\selectfont,
    text=galkInk
  },
  galk/small label/.style={
    font=\rmfamily\fontsize{8.7pt}{9.7pt}\selectfont,
    text=galkInk
  },
  galk/zone label/.style={
    font=\rmfamily\bfseries\fontsize{8.7pt}{9.7pt}\selectfont,
    text=galkInk,
    anchor=north west,
    inner sep=.5pt
  },
  galk/entanglement zone/.style={
    draw=galkEntangle,
    fill=galkEntangleFill,
    line width=.40pt
  },
  galk/storage zone/.style={
    draw=galkStorage,
    fill=galkStorageFill,
    line width=.40pt
  },
  galk/site/.style={
    circle,
    draw=galkInk!68,
    fill=white,
    line width=.36pt,
    minimum size=1.35mm,
    inner sep=0pt
  },
  galk/occupied atom/.style={
    circle,
    draw=galkInk,
    fill=galkInk,
    line width=.36pt,
    minimum size=1.55mm,
    inner sep=0pt
  },
  galk/target site/.style={
    circle,
    draw=galkInk!60,
    fill=galkInk!10,
    line width=.34pt,
    minimum size=1.55mm,
    inner sep=0pt
  },
  galk/q label/.style={
    font=\rmfamily\itshape\fontsize{8.5pt}{9.5pt}\selectfont,
    text=galkInk,
    inner sep=.25pt
  },
  galk/rydberg pair/.style={
    draw=galkInk!62,
    densely dashed,
    line width=.36pt
  },
  galk/accepted path/.style={
    -{Stealth[length=.95mm,width=.62mm]},
    draw=galkInk,
    line width=.48pt
  },
  galk/alternative path/.style={
    -{Stealth[length=.95mm,width=.62mm]},
    draw=galkMuted,
    densely dashed,
    line width=.48pt
  },
  galk/future path/.style={
    -{Stealth[length=.95mm,width=.62mm]},
    draw=galkLook,
    densely dashed,
    line width=.48pt
  },
  galk/invalid path/.style={
    -{Stealth[length=.95mm,width=.62mm]},
    draw=galkBad,
    densely dashed,
    line width=.50pt
  }
}
\fi

\begin{tikzpicture}[
  galk/base,x=1cm,y=1cm,
  trap/.style={galk/site},
  atom/.style={galk/occupied atom},
  transition/.style={galk/accepted path},
  panel title/.style={galk/panel title},
  state note/.style={font=\rmfamily\fontsize{8.7pt}{9.7pt}\selectfont},
  atom label/.style={galk/q label,font=\rmfamily\itshape\fontsize{8.8pt}{9.8pt}\selectfont},
  pulse/.style={draw=galkEntangle,densely dashed,line width=.42pt}
]
  \path[use as bounding box] (0,0) rectangle (17.80,5.85);

  % One regular optical-trap geometry is reused without deformation.
  \def\miniDevice#1#2{%
    \begin{scope}[shift={(#1,#2)}]
      \path[galk/entanglement zone] (0,.80) rectangle (1.12,1.30);
      \path[galk/storage zone] (0,0) rectangle (1.12,.70);
      \draw[galk/rydberg pair] (.27,1.05) ellipse (.15 and .13);
      \draw[galk/rydberg pair] (.85,1.05) ellipse (.15 and .13);
      \foreach \xx in {.18,.36,.76,.94}
        \node[trap] at (\xx,1.05) {};
      \foreach \xx in {.18,.56,.94} {
        \node[trap] at (\xx,.20) {};
        \node[trap] at (\xx,.50) {};
      }
    \end{scope}%
  }
  \def\gatePair#1#2#3#4{%
    \node[atom] at (#1,#2) {};
    \node[atom] at (#1+.18,#2) {};
    \draw[galkInk,line width=.42pt] (#1,#2)--(#1+.18,#2);
    \node[atom label,anchor=south] at (#1+.09,#2+.30) {$#3,#4$};
  }
  \def\siteAnnotation#1#2#3#4{%
    \node[state note] at (#1+.56,#2-.27) {$q$ at $#3$};
    \draw[galkInk!65,line width=.35pt]
      (#1+.38,#2-.09)--(#1+.38,#2+#4-.16)--(#1+.49,#2+#4-.08);
  }

  \draw[galkInk!24,line width=.38pt] (5.48,.35)--(5.48,5.58);
  \draw[galkInk!24,line width=.38pt] (11.15,.35)--(11.15,5.58);

  % (a) The interaction order determines when q is needed again.
  \node[panel title] at (.15,5.57) {(a) Non-adjacent interaction};
  \foreach \xx/\lab in {.20/$L_{\ell}$,1.49/$L_{\ell+1}$,
                       2.78/$L_{\ell+2}$,4.07/$L_{\ell+3}$} {
    \miniDevice{\xx}{3.08}
    \node[anchor=south] at (\xx+.56,4.92) {\lab};
  }
  \gatePair{.38}{4.13}{q}{p}
  \gatePair{2.25}{4.13}{u}{v}
  \gatePair{3.54}{4.13}{w}{x}
  \gatePair{4.25}{4.13}{q}{r}
  \foreach \xx in {1.67,2.96} {
    \node[atom] at (\xx,3.58) {};
    \node[atom label,anchor=center] at (\xx+.17,3.40) {$q$};
  }
  \node[state note,anchor=north] at (.76,2.88) {$\mathrm{CZ}(q,p)$};
  \node[state note,anchor=north] at (2.05,2.88) {$\mathrm{CZ}(u,v)$};
  \node[state note,anchor=north] at (3.34,2.88) {$\mathrm{CZ}(w,x)$};
  \node[state note,anchor=north] at (4.63,2.88) {$\mathrm{CZ}(q,r)$};
  \draw[galkInk!60,{Stealth[length=.75mm]}-{Stealth[length=.75mm]},line width=.40pt]
    (1.72,2.18)--(4.09,2.18);
  \node[fill=white,inner sep=1pt,state note] at (2.91,2.18) {two intervening layers};

  % The two-line zone key uses the same rectangular bands as the snapshots.
  \path[galk/entanglement zone] (.45,1.13) rectangle (.73,1.34);
  \node[anchor=west,state note] at (.88,1.235) {entanglement zone};
  \path[galk/storage zone] (.45,.60) rectangle (.73,.81);
  \node[anchor=west,state note] at (.88,.705) {storage zone};

  % (b) The gate sequence is identical; only q's inter-layer location differs.
  \node[panel title] at (5.68,5.57) {(b) Inter-layer residency};
  % Layer headers share the same top band as panels (a) and (c).
  % Place each residency description below its own row of snapshots.
  \node[anchor=west] at (5.73,2.80) {Entanglement-resident};
  \node[anchor=west] at (5.73,.53) {Storage-mediated};
  \foreach \xx/\lab in {5.73/$L_{\ell}$,7.04/$L_{\ell+1}$,
                       8.35/$L_{\ell+2}$,9.66/$L_{\ell+3}$} {
    \miniDevice{\xx}{3.08}
    \miniDevice{\xx}{.87}
    \node[anchor=south] at (\xx+.56,4.92) {\lab};
  }
  \gatePair{5.91}{4.13}{q}{p}
  \gatePair{7.80}{4.13}{u}{v}
  \gatePair{9.11}{4.13}{w}{x}
  \gatePair{9.84}{4.13}{q}{r}
  \foreach \xx in {7.22,8.53} {
    \node[atom] at (\xx,4.13) {};
    \node[atom label,anchor=south] at (\xx,4.43) {$q$};
    \draw[pulse] (\xx-.13,3.94)--(\xx-.13,4.32);
    \draw[pulse] (\xx+.13,3.94)--(\xx+.13,4.32);
  }
  \gatePair{5.91}{1.92}{q}{p}
  \gatePair{7.80}{1.92}{u}{v}
  \gatePair{9.11}{1.92}{w}{x}
  \gatePair{9.84}{1.92}{q}{r}
  \foreach \xx in {7.22,8.53} {
    \node[atom] at (\xx,1.37) {};
    \node[atom label,anchor=center] at (\xx+.17,1.19) {$q$};
  }

  % (c) s1 and s2 are real traps on the same regular storage lattice.
  \node[panel title] at (11.35,5.57) {(c) Storage-site dependence};
  \foreach \xx/\lab in {11.48/$L_{\ell}$,13.84/$L_{\ell+1}$,16.20/$L_{\ell+2}$} {
    \miniDevice{\xx}{3.08}
    \miniDevice{\xx}{.94}
    \node[anchor=south] at (\xx+.56,4.92) {\lab};
  }

  % The source state contains q in the entanglement zone and u in storage.
  \foreach \yy in {3.08,.94} {
    \node[atom] at (11.66,\yy+1.05) {};
    \node[atom label,anchor=south] at (11.66,\yy+1.35) {$q$};
    \node[atom] at (11.66,\yy+.50) {};
    \node[atom label,anchor=center] at (11.83,\yy+.32) {$u$};
  }
  \draw[transition] (12.78,3.66)--(13.66,3.66);
  \node[state note,anchor=south] at (13.22,3.80) {$q\to s_1$};
  \draw[transition] (12.78,1.52)--(13.66,1.52);
  \node[state note,anchor=south] at (13.22,1.66) {$q\to s_2$};

  % At s1, q occupies the middle trap of u's direct route.
  \foreach \xx in {13.84,16.20} {
    \node[atom] at (\xx+.18,3.58) {};
    \node[atom] at (\xx+.56,3.58) {};
    \draw[galk/invalid path] (\xx+.26,3.58)--(\xx+.86,3.58);
    \draw[galkBad,line width=.60pt] (\xx+.48,3.50)--(\xx+.64,3.66);
    \draw[galkBad,line width=.60pt] (\xx+.48,3.66)--(\xx+.64,3.50);
  }
  \siteAnnotation{13.84}{3.08}{s_1}{.50}
  \siteAnnotation{16.20}{3.08}{s_1}{.50}
  \draw[galk/invalid path] (15.14,3.66)--(16.02,3.66);
  \node[state note,anchor=south,text=galkBad] at (15.58,3.80) {Blocked};

  % At s2, the same single-atom route remains clear and u reaches its target.
  \node[atom] at (14.02,1.44) {};
  \node[atom] at (14.40,1.14) {};
  \draw[transition] (14.10,1.44)--(14.70,1.44);
  \node[atom] at (17.14,1.44) {};
  \node[atom] at (16.76,1.14) {};
  \node[atom label,anchor=center] at (16.97,1.26) {$u$};
  \siteAnnotation{13.84}{.94}{s_2}{.20}
  \siteAnnotation{16.20}{.94}{s_2}{.20}
  \draw[transition] (15.14,1.52)--(16.02,1.52);
  \node[state note,anchor=south] at (15.58,1.66) {Clear};
\end{tikzpicture}
}
  \caption{非相邻交互下的驻留与存储位选择。(a) $q$在$L_\ell$与$L_{\ell+3}$参与交互。(b) 驻留承担中间层受激损失，回迁承担往返输运损失；橙色虚线标示中间两层对驻留原子的Rydberg脉冲，仅保留$q$及相关门对。(c) $s_1$、$s_2$为实际存储光阱，$q$的回迁位置决定$u$的后续直达路径是否可行。各状态共用相同区域几何，尺寸与阱位数量为示意。}
  \label{fig:baseline-motivation}
\end{figure*}

在原子位置基本保持固定的阵列上，编译器通过电路分块和初始布局，使电路匹配器件能够执行的门类型及其静态连接。Patel等人\cite{patel2022geyser}的Geyser围绕原生多比特门划分电路，GRAPHINE\cite{patel2023graphine}依据带权交互图选择原子在阵列中的位置。Schmid等人\cite{schmid2024hybrid}结合长程门与原子输运，以适配同时具有多种连接方式的中性原子器件。

不区分存储区与纠缠区的可重构阵列主要通过原子移动改善连通性。OLSQ-DPQA联合求解门调度、放置与移动，Enola将双比特门调度和并行搬运分解处理\cite{tan2022mapping,tan2024olsqdpqa,tan2025enola}。Atomique与PARALLAX分别利用多AOD阵列和无SWAP移动扩大并行性\cite{wang2024atomique,ludmir2024parallax}。Q-Pilot将数据原子固定于静态阱，通过可回收的移动辅助原子完成双比特交互；Weaver通过专用中间表示支持不同可重构阵列之间的编译适配\cite{wang2024qpilot,kirmemis2025weaver}。这些方法关注门交互的实现与并行执行，分区架构则还需考虑存储隔离和跨区输运。

面向分区架构，Lin等人\cite{lin2025zac}提出复用感知（Reuse-Aware）编译器ZAC，以相邻层匹配保留可复用的纠缠区原子，并在门位与存储位匹配中计入后续双比特门的配对关系及其距离代价。Stade等人\cite{stade2025routingaware}在AOD抽象模型\cite{stade2024abstractmodel}上提出路由感知（Routing-Aware）放置，用A*搜索门放置和中间放置，代价包含相容移动分组、组内最长轨迹，以及后续双比特门配对原子的间距。前者着重减少重复阱转移，后者使位置选择与路由并行性相联系。Huang等人\cite{huang2026zap}的ZAP同样面向分区架构：层序加权初始化并估计早期移动冲突，动态放置先按下次使用前的受激损失决定驻留或回迁，再逐门选择目标位置，以距离和冲突数作顺序贪心估价。ZAP公布的阵列间距、相干损失模型和保真度统计项与本文设置不同，故不将其原文数值纳入直接比较；实验在统一输入、架构与评分模型下比较GA-LK、复用感知编译和路由感知放置。

\subsection{跨层驻留与存储位选择}

跨层驻留决定原子在两次交互之间所在的区域。留在纠缠区可减少往返转移，回到存储区可避免中间层的Rydberg暴露。回迁存储位决定后续搬运的起点，以及是否占用其他原子的通道。驻留、回迁存储位与后续门位共同决定了层间重排的物理代价。

原子是否值得驻留，取决于它何时再次参与交互。原子$q$在$L_\ell$与$p$执行CZ门，到$L_{\ell+3}$才与$r$交互，中间两层的门均不包含$q$，见图~\ref{fig:baseline-motivation}(a)。若$q$留在纠缠区，就会经历这两层的空闲受激；若回到存储区，则需要为下一次交互再进入纠缠区，增加阱转移和输运，见图~\ref{fig:baseline-motivation}(b)。因此，驻留选择需要同时考虑再次使用的层距和实际往返代价。

回迁位置同时会改变其他原子的可行路径。图~\ref{fig:baseline-motivation}(c)中的$s_1$、$s_2$是存储阵列内两个不同的静态光阱，空心圆表示原子$u$在下一层的目标位置。$q$回到$s_1$后，会阻挡$u$的直达路径；选择通道之外的$s_2$则可保留该路径。两种选择具有相同的当前阱转移数，但前者需改变后续路径或增加中转。按未来层序更新原子位置和阱位占用，能够在选择回迁位置时计入这种差异。

\subsection{本文贡献}
\label{sec:contributions}

本文提出结合遗传算法与多层前瞻的联合放置方法GA-LK。
\begin{itemize}
  \item \PaperRevision{门位、驻留与回迁存储位相互制约。本文将这些耦合选择建模为保真度损失最小化问题，并以遗传搜索与匹配求解。遗传搜索联合选择双比特门位和原子的驻留或回迁方案，匹配算法为回迁原子分配互不重复的存储阱位；每组完整方案均按阱转移、空闲受激和相干损失评价。}
  \item \PaperRevision{当前放置会改变后续层的输运与受激条件。本文按各候选完成当前层后的原子位置和阱位占用，逐层估计后续执行损失，再按层距衰减其权重，使驻留和存储位选择兼顾多层执行代价。这一评价还用于默认初始化：比较不同起始存储布局下早期门层的执行损失，在电路开始前选择初始位置。}
\end{itemize}
